\documentclass[aps,prb,preprint,amsmath,amssymb,superscriptaddress,longbibliography,floatfix]{revtex4-2}

\usepackage[T1]{fontenc}
\usepackage[utf8]{inputenc}
\usepackage{graphicx}
\usepackage{xcolor}
\usepackage{booktabs}
\usepackage{makecell}
\definecolor{TDKgreen}{rgb}{0,0.45,0}
\newcommand{\TDKin}[1]{\textcolor{blue}{#1}}
\newcommand{\TDKout}[1]{\textcolor{red}{#1}}
\newcommand{\TDKcomment}[1]{\textcolor{TDKgreen}{\textbf{[TDK:} #1\textbf{]}}}

\begin{document}

\title{Supplemental Material for ``The UZH protocol: Separating and reducing Gaussian-basis and pseudopotential errors in CP2K''}

\author{Hossein Mirhosseini}
\affiliation{Center for Advanced Systems Understanding (CASUS), Conrad-Schiedt-Stra{\ss}e 20, 02826 G\"orlitz, Germany}
\affiliation{Helmholtz-Zentrum Dresden-Rossendorf, Bautzner Landstra{\ss}e 400, 01328 Dresden, Germany}
\author{Tiziano M. A. M\"uller}
\affiliation{Department of Chemistry, University of Zurich, CH-8057 Zurich, Switzerland}
\author{Matthias Krack}
\affiliation{Paul Scherrer Institute, 5232 Villigen PSI, Switzerland}
\author{Thomas D. K\"uhne}
\email{tkuehne@cp2k.org}
\affiliation{Center for Advanced Systems Understanding (CASUS), Conrad-Schiedt-Stra{\ss}e 20, 02826 G\"orlitz, Germany}
\affiliation{Helmholtz-Zentrum Dresden-Rossendorf, Bautzner Landstra{\ss}e 400, 01328 Dresden, Germany}
\affiliation{Institute of Artificial Intelligence, Technische Universit\"at Dresden, Helmholtzstra{\ss}e 10, 01069 Dresden, Germany}
\author{J\"urg Hutter}
\affiliation{Department of Chemistry, University of Zurich, CH-8057 Zurich, Switzerland}

\maketitle

\section{Source material and implementation details}

The main text summarizes the protocol, while this Supplemental Material records implementation, calibration, and diagnostic details needed to reproduce and extend the UZH settings. The organization follows the original molecularly optimized (MOLOPT), Goedecker--Teter--Hutter (GTH), and Hartwigsen--Goedecker--Hutter (HGH) literature~\cite{VandeVondele2007,Goedecker1996,Hartwigsen1998}, the revised MOLOPT benchmarking work~\cite{MullerThesis2024}, the CP2K implementation papers~\cite{CP2K2020,CP2KMadeSimple2026}, and the Automated Interactive Infrastructure and Database for Computational Science (AiiDA) verification protocol and its Supplemental Material~\cite{Huber2021,Bosoni2024}.

\subsection{CP2K ATOM and basis optimization}

The updated CP2K ATOM implementation developed for the present UZH work is exposed through a dedicated input section for atomic calculations. The relevant run types are energy calculations, basis optimization, and pseudopotential optimization. The ATOM method section supports restricted and unrestricted Kohn-Sham and Hartree-Fock atomic calculations and scalar-relativistic Hamiltonians including zeroth-order regular approximation (ZORA)-type and Douglas--Kroll--Hess variants. For the UZH pseudopotential fits reported here, the all-electron atomic references were generated with the third-order Douglas--Kroll--Hess (DKH3) Hamiltonian\TDKin{~\cite{DouglasKroll1974,Hess1986,NakajimaHirao2000}}. The relativistic section allows the order of the Douglas--Kroll--Hess transformation, the transformation strategy, and the ZORA model-potential variant to be selected. The potential section supports all-electron, GTH, unified pseudopotential format, separable Gaussian pseudopotential, and effective core potential inputs. For the GTH format, the input layer exposes a CP2K-standard GTH potential block, and the print layer can write the optimized potential in CP2K format and in a form compatible with the unified pseudopotential format.

The CP2K basis optimizer works from template, work, and output basis files. It supports simultaneous fitting of related basis levels and allows the condition number of the overlap matrix to enter the objective. This mirrors the original MOLOPT design: diffuse flexibility is desirable for weak interactions and response properties, but the final basis must remain numerically stable in condensed-phase simulations.

The pseudopotential optimizer constructs both all-electron and pseudopotential atomic references. Occupations are partitioned into core and valence contributions; semicore choices are therefore explicit in the input and can be changed when the verification workflow reveals poor transferability. The optimizer evaluates the all-electron reference with the DKH3 Hamiltonian used here, constructs the corresponding pseudopotential atom, and fits the compact GTH parameters with a derivative-free Powell algorithm. The fitted target is therefore not an isolated numerical object but a DKH3-based scalar-relativistic atomic mapping from an all-electron reference to the analytic separable GTH form. This machinery is what allows the UZH protocol to improve the pseudopotential component once the comparison between SIRIUS-GTH-UZH and the SIRIUS all-electron reference identifies a pseudopotential-limited element.

\subsection{SIRIUS calculations}

SIRIUS is used in two complementary modes~\cite{Zhang2022Computation,Zhang2023Sirius,SiriusDocs}. In pseudopotential plane-wave mode, SIRIUS-GTH-UZH uses the same GTH-UZH pseudopotentials as CP2K but removes the localized Gaussian-basis approximation. In all-electron full-potential linearized augmented-plane-wave (FP-LAPW) mode, SIRIUS provides the reference equation of state. The CP2K-GTH-UZH versus SIRIUS-GTH-UZH comparison isolates the CP2K Gaussian-basis error, while the two SIRIUS modes isolate the pseudopotential error.

In the FP-LAPW calculation the unit cell is divided into non-overlapping muffin-tin spheres and an interstitial region. Radial functions multiplied by spherical harmonics describe the wave functions inside the spheres, while plane waves describe the interstitial region; local orbitals are used to improve the description of semicore and high-lying states. Because the method is full-potential, the charge density and Kohn-Sham potential are represented without imposing a muffin-tin shape approximation. This all-electron reference is therefore well suited for the present purpose: any difference between SIRIUS-GTH-UZH and SIRIUS-FP-LAPW primarily measures the frozen-core and pseudization error of the GTH-UZH potential, including the quality of semicore and scalar-relativistic transferability.

For heavy elements the relativistic settings are part of the definition of the reference. The SIRIUS full-potential species files contain the radial grid, linearized augmented-plane-wave (LAPW)/augmented-plane-wave plus local-orbital (APW+lo) basis information, core-state treatment, and valence-relativity choice. The FP-LAPW reference calculations use the infinite-order regular approximation (IORA) for valence scalar relativity~\cite{Dyall1999IORA}; this IORA setting is recorded in the full-potential species definition and kept fixed throughout the equation-of-state comparison. The CP2K interface exposes the electronic-structure mode as either pseudopotential or full-potential LAPW+lo, and the full-potential setup distinguishes core relativity from valence relativity. In the UZH protocol we compare scalar-relativistic pseudopotentials, fitted to DKH3 atomic references in CP2K ATOM, to scalar-relativistic all-electron FP-LAPW references; explicit spin-orbit splittings are outside the present benchmark. This convention keeps the error decomposition aligned with the scalar-relativistic GTH/HGH pseudopotentials used by CP2K/Quickstep.

\section{Molecular calibration}

The molecular calibration assesses the revised UZH-MOLOPT basis sequence before the solid-state verification step. The benchmark contains small molecules covering main-group bonds, polar bonds, nonlinear and dihedral angles, dipole moments, polarizabilities, and vibrational frequencies. The underlying CP2K molecular test calculations are provided in the MolTest Zenodo archive, which contains calculations with revised CP2K basis sets and pseudopotentials and CP2K all-electron basis sets for the Perdew--Burke--Ernzerhof (PBE), Perdew--Burke--Ernzerhof hybrid (PBE0), and Tao--Perdew--Staroverov--Scuseria (TPSS) exchange-correlation functionals~\cite{Perdew1996,Adamo1999,Tao2003TPSS,HutterMuller2023MolTest}. The molecular basis sequence follows the main text and includes double-zeta valence plus polarization (DZVP), triple-zeta valence plus polarization (TZVP), and triple-zeta valence plus two polarization shells (TZV2P) levels; the CP2K all-electron reference calculations use the Gaussian-and-augmented-plane-wave (GAPW) method. For the reference-code comparison below, Gaussian~16 uses the def2 quadruple-zeta valence plus polarization (def2-QZVP) basis, whereas CP2K/GAPW uses a quadruple-zeta valence plus two polarization shells (QZVPP) Gaussian basis. Figures~\ref{fig:si_molecular_angles}, \ref{fig:si_molecular_dihedrals}, and \ref{fig:si_molecular_dipoles} provide the angular and dipole-moment statistics that complement the main-text bond-length and polarizability analysis. Figures~\ref{fig:si_bond_classes} and \ref{fig:si_reference_comparison} summarize two additional molecular checks from the thesis-level calibration data: the bond-class dependence of the structural errors and the consistency of Gaussian~16 and CP2K all-electron reference calculations. The role of this stage is to ensure chemically balanced transferability and numerical stability before the basis is used in periodic calculations.

\begin{figure}
\centering
\includegraphics[width=0.9\linewidth]{TM_mol_angles.png}
\caption{\label{fig:si_molecular_angles}
Distribution of relative errors for nonlinear angles in the molecular calibration set. Angular observables are more sensitive than bond lengths to automated atom matching and near-symmetric geometries, but the systematic improvement from DZVP to TZV2P remains visible.}
\end{figure}

\begin{figure}
\centering
\includegraphics[width=0.9\linewidth]{TM_mol_dihedrals.png}
\caption{\label{fig:si_molecular_dihedrals}
Distribution of relative errors for dihedral angles in the molecular calibration set. The broader tails reflect the stronger sensitivity of torsional coordinates to small structural and matching differences, while the basis-set hierarchy remains consistent with the bond-length and nonlinear-angle tests.}
\end{figure}

\begin{figure}
\centering
\includegraphics[width=0.9\linewidth]{TM_mol_dipoles.png}
\caption{\label{fig:si_molecular_dipoles}
Distribution of relative errors for molecular dipole moments, adapted from Fig.~4.11 of Ref.~\cite{MullerThesis2024}. The statistics include converged structures with non-vanishing dipole moments and show that response-sensitive density properties retain the same systematic basis-set improvement observed for structural observables.}
\end{figure}

\begin{figure}[p]
\includegraphics[width=\linewidth,height=0.82\textheight,keepaspectratio]{TM_thesis_fig_4_13.png}
\caption{\label{fig:si_bond_classes}
Bond-class-resolved distribution of relative bond-length errors across the molecular calibration set, adapted from Fig.~4.13 of Ref.~\cite{MullerThesis2024}. Outliers are hidden and the C--C panel uses a smaller vertical scale. The grouping shows that the UZH-MOLOPT basis sequence improves systematically across chemically distinct bond classes before periodic verification.}
\end{figure}

\begin{figure}[p]
\includegraphics[width=\linewidth,height=0.82\textheight,keepaspectratio]{TM_thesis_fig_4_16.png}
\caption{\label{fig:si_reference_comparison}
Comparison of molecular calibration errors obtained with Gaussian~16/def2-QZVP and CP2K/GAPW-QZVPP all-electron references, adapted from Fig.~4.16 of Ref.~\cite{MullerThesis2024}. Relative errors are shown for bond lengths, dipole moments, and average molecular polarizabilities, while absolute errors are shown for angles and dihedral angles. The close agreement supports the use of the CP2K all-electron molecular reference in the UZH calibration loop.}
\end{figure}

\clearpage

\section{Additional improved-setting diagnostics}

Figure~\ref{fig:si_pp_residual_update} provides two complementary diagnostics of the pseudopotential part of the UZH optimization. The upper panel measures the remaining pseudopotential transferability error of the improved SIRIUS-GTH-UZH-new settings against the SIRIUS-FP-LAPW all-electron reference. The lower panel compares SIRIUS-GTH-UZH-new directly with the original SIRIUS-GTH-UZH setup and therefore shows where the pseudopotential update changes the equation-of-state curves. Large changes in the lower panel that coincide with small residual errors in the upper panel identify successful refits, whereas elements that remain large in the upper panel mark residual pseudopotential errors not removed by the present update.

\begin{figure}
\centering
\includegraphics[width=0.9\linewidth]{periodic-table-unaries-pp-residual-and-update.png}
\caption{\label{fig:si_pp_residual_update}
Pseudopotential residual and update diagnostics for the improved UZH settings. Top: remaining pseudopotential contribution of SIRIUS-GTH-UZH-new relative to the SIRIUS-FP-LAPW all-electron reference. Bottom: direct comparison of SIRIUS-GTH-UZH-new with the original SIRIUS-GTH-UZH setup, showing where the pseudopotential update changes the equation-of-state curves.}
\end{figure}

\clearpage

\section{Representative equation-of-state curves}

The periodic-table heat maps in the main text are compact summaries of many equation-of-state calculations. Figures~\ref{fig:si_cr}, \ref{fig:si_ba}, \ref{fig:si_ne}, \ref{fig:si_au}, and \ref{fig:si_ga} show representative raw equation-of-state comparisons for elements that illustrate different error regimes. They are included to connect the color-coded $\varepsilon$ maps to the underlying energy--volume curves: Cr illustrates a pseudopotential-limited case, Ba and Ne emphasize Gaussian-basis sensitivity, and Au and Ga represent heavier elements for which basis completeness and scalar-relativistic pseudopotential transferability must both be monitored.

\begin{figure}
\includegraphics[width=0.9\linewidth]{Cr.pdf}
\caption{\label{fig:si_cr}
Representative equation-of-state curves for Cr. The behavior is characteristic of a pseudopotential-limited case in the original settings.}
\end{figure}

\begin{figure}
\includegraphics[width=0.9\linewidth]{Ba.pdf}
\caption{\label{fig:si_ba}
Representative equation-of-state curves for Ba. The comparison highlights a case with a substantial Gaussian-basis contribution in the original CP2K-GTH-UZH settings.}
\end{figure}

\begin{figure}
\includegraphics[width=0.9\linewidth]{Ne.pdf}
\caption{\label{fig:si_ne}
Representative equation-of-state curves for Ne. Noble gases are sensitive tests of diffuse basis flexibility because their equation-of-state curves are shallow.}
\end{figure}

\begin{figure}
\includegraphics[width=0.9\linewidth]{Au.pdf}
\caption{\label{fig:si_au}
Representative equation-of-state curves for Au, illustrating the behavior of a heavy element where basis completeness and relativistic pseudopotential transferability must both be monitored.}
\end{figure}

\begin{figure}
\includegraphics[width=0.9\linewidth]{Ga.pdf}
\caption{\label{fig:si_ga}
Representative equation-of-state curves for Ga.}
\end{figure}

\section{Release repository organization}

The release repository for the UZH protocol is hosted in the DCM-Uni-Paderborn GitHub organization under the repository name UZH-protocol. It is organized as follows:
\begin{enumerate}
\item CP2K-GTH-UZH AiiDA workflow inputs and parsed results.
\item SIRIUS-GTH-UZH AiiDA workflow inputs and parsed results.
\item SIRIUS-FP-LAPW all-electron workflow inputs and parsed results.
\item Optimized UZH-MOLOPT basis files and GTH-UZH pseudopotential files.
\item Scripts used to evaluate the $\varepsilon$ metric and generate the periodic-table maps.
\item A manifest documenting code versions, basis-file hashes, pseudopotential hashes, and workflow protocol names.
\end{enumerate}

\bibliography{bib_prb_tdk}

\end{document}
